\IfFileExists{stacks-project.cls}{%
\documentclass{stacks-project}
}{%
\documentclass{amsart}
}
\usepackage{amssymb}

% For dealing with references we use the comment environment
\usepackage{verbatim}
\newenvironment{reference}{\comment}{\endcomment}
%\newenvironment{reference}{}{}
\newenvironment{slogan}{\comment}{\endcomment}
\newenvironment{history}{\comment}{\endcomment}

% For commutative diagrams we use Xy-pic
\usepackage[all]{xy}

% We use 2cell for 2-commutative diagrams.
\xyoption{2cell}
\UseAllTwocells

% We use multicol for the list of chapters between chapters
\usepackage{multicol}

% This is generally recommended for better output
\usepackage{lmodern}
\usepackage[T1]{fontenc}
\usepackage{textalpha}

% For cross-file-references
% We prefer the snapshot tmp/xref/#2.aux,
% which latexmkrc replaces atomically, so
% that documents compiled simultaneously
% never read a half-written .aux file.
% The last argument keeps links to #2.pdf.
\usepackage{xr-hyper}
\newcommand{\maybeexternaldocument}[2]{%
  \IfFileExists{#2.aux}{%
    \externaldocument[#1]{#2}%
  }{%
    \PackageWarning{stack}{Missing #2.aux; external references with prefix #1 unavailable}%
  }%
}
\let\externaldocumentorig\externaldocument
\renewcommand{\externaldocument}[2][]{%
  \IfFileExists{tmp/xref/#2.aux}{%
    \externaldocumentorig[#1]{tmp/xref/#2}[#2.pdf]%
  }{%
  \IfFileExists{#2.aux}{%
    \externaldocumentorig[#1]{#2}%
  }{%
    \PackageWarning{stack}{Missing #2.aux; external references with prefix #1 unavailable}%
  }%
  }%
}

% Package for hypertext links:
\usepackage{hyperref}

% For any local file, say "hello.tex" you want to link to please
% use \externaldocument[hello-]{hello}
\externaldocument[introduction-]{introduction}
\externaldocument[conventions-]{conventions}
\externaldocument[sets-]{sets}
\externaldocument[categories-]{categories}
\externaldocument[topology-]{topology}
\externaldocument[sheaves-]{sheaves}
\externaldocument[sites-]{sites}
\externaldocument[stacks-]{stacks}
\externaldocument[fields-]{fields}
\externaldocument[algebra-]{algebra}
\externaldocument[brauer-]{brauer}
\externaldocument[homology-]{homology}
\externaldocument[derived-]{derived}
\externaldocument[simplicial-]{simplicial}
\externaldocument[more-algebra-]{more-algebra}
\externaldocument[smoothing-]{smoothing}
\externaldocument[modules-]{modules}
\externaldocument[sites-modules-]{sites-modules}
\externaldocument[injectives-]{injectives}
\externaldocument[cohomology-]{cohomology}
\externaldocument[sites-cohomology-]{sites-cohomology}
\externaldocument[dga-]{dga}
\externaldocument[dpa-]{dpa}
\externaldocument[sdga-]{sdga}
\externaldocument[hypercovering-]{hypercovering}
\externaldocument[schemes-]{schemes}
\externaldocument[constructions-]{constructions}
\externaldocument[properties-]{properties}
\externaldocument[morphisms-]{morphisms}
\externaldocument[coherent-]{coherent}
\externaldocument[divisors-]{divisors}
\externaldocument[limits-]{limits}
\externaldocument[varieties-]{varieties}
\externaldocument[topologies-]{topologies}
\externaldocument[descent-]{descent}
\externaldocument[perfect-]{perfect}
\externaldocument[more-morphisms-]{more-morphisms}
\externaldocument[flat-]{flat}
\externaldocument[groupoids-]{groupoids}
\externaldocument[more-groupoids-]{more-groupoids}
\externaldocument[etale-]{etale}
\externaldocument[chow-]{chow}
\externaldocument[intersection-]{intersection}
\externaldocument[pic-]{pic}
\externaldocument[weil-]{weil}
\externaldocument[adequate-]{adequate}
\externaldocument[dualizing-]{dualizing}
\externaldocument[duality-]{duality}
\externaldocument[discriminant-]{discriminant}
\externaldocument[derham-]{derham}
\externaldocument[local-cohomology-]{local-cohomology}
\externaldocument[algebraization-]{algebraization}
\externaldocument[curves-]{curves}
\externaldocument[resolve-]{resolve}
\externaldocument[models-]{models}
\externaldocument[functors-]{functors}
\externaldocument[equiv-]{equiv}
\externaldocument[pione-]{pione}
\externaldocument[etale-cohomology-]{etale-cohomology}
\externaldocument[proetale-]{proetale}
\externaldocument[relative-cycles-]{relative-cycles}
\externaldocument[more-etale-]{more-etale}
\externaldocument[trace-]{trace}
\externaldocument[crystalline-]{crystalline}
\externaldocument[spaces-]{spaces}
\externaldocument[spaces-properties-]{spaces-properties}
\externaldocument[spaces-morphisms-]{spaces-morphisms}
\externaldocument[decent-spaces-]{decent-spaces}
\externaldocument[spaces-cohomology-]{spaces-cohomology}
\externaldocument[spaces-limits-]{spaces-limits}
\externaldocument[spaces-divisors-]{spaces-divisors}
\externaldocument[spaces-over-fields-]{spaces-over-fields}
\externaldocument[spaces-topologies-]{spaces-topologies}
\externaldocument[spaces-descent-]{spaces-descent}
\externaldocument[spaces-perfect-]{spaces-perfect}
\externaldocument[spaces-more-morphisms-]{spaces-more-morphisms}
\externaldocument[spaces-flat-]{spaces-flat}
\externaldocument[spaces-groupoids-]{spaces-groupoids}
\externaldocument[spaces-more-groupoids-]{spaces-more-groupoids}
\externaldocument[bootstrap-]{bootstrap}
\externaldocument[spaces-pushouts-]{spaces-pushouts}
\externaldocument[spaces-chow-]{spaces-chow}
\externaldocument[groupoids-quotients-]{groupoids-quotients}
\externaldocument[spaces-more-cohomology-]{spaces-more-cohomology}
\externaldocument[spaces-simplicial-]{spaces-simplicial}
\externaldocument[spaces-duality-]{spaces-duality}
\externaldocument[formal-spaces-]{formal-spaces}
\externaldocument[restricted-]{restricted}
\externaldocument[spaces-resolve-]{spaces-resolve}
\externaldocument[formal-defos-]{formal-defos}
\externaldocument[defos-]{defos}
\externaldocument[cotangent-]{cotangent}
\externaldocument[examples-defos-]{examples-defos}
\externaldocument[algebraic-]{algebraic}
\externaldocument[examples-stacks-]{examples-stacks}
\externaldocument[stacks-sheaves-]{stacks-sheaves}
\externaldocument[criteria-]{criteria}
\externaldocument[artin-]{artin}
\externaldocument[quot-]{quot}
\externaldocument[stacks-properties-]{stacks-properties}
\externaldocument[stacks-morphisms-]{stacks-morphisms}
\externaldocument[stacks-limits-]{stacks-limits}
\externaldocument[stacks-cohomology-]{stacks-cohomology}
\externaldocument[stacks-perfect-]{stacks-perfect}
\externaldocument[stacks-introduction-]{stacks-introduction}
\externaldocument[stacks-more-morphisms-]{stacks-more-morphisms}
\externaldocument[stacks-geometry-]{stacks-geometry}
\externaldocument[moduli-]{moduli}
\externaldocument[moduli-curves-]{moduli-curves}
\externaldocument[examples-]{examples}
\externaldocument[exercises-]{exercises}
\externaldocument[guide-]{guide}
\externaldocument[desirables-]{desirables}
\externaldocument[coding-]{coding}
\externaldocument[obsolete-]{obsolete}
\externaldocument[fdl-]{fdl}
\externaldocument[index-]{index}

% Heap Project documents
\externaldocument[linear-algebra-]{linear-algebra}
\externaldocument[groups-]{groups}
\externaldocument[complex-analysis-]{complex-analysis}
\externaldocument[functional-analysis-]{functional-analysis}
\externaldocument[categories-]{categories}
\externaldocument[commutative-algebra-]{commutative-algebra}
\externaldocument[ringed-spaces-]{ringed-spaces}
\externaldocument[homological-algebra-]{homological-algebra}
\externaldocument[lie-algebras-]{lie-algebras}
\externaldocument[representation-theory-]{representation-theory}
\externaldocument[smooth-manifolds-]{smooth-manifolds}
\externaldocument[differential-topology-]{differential-topology}
\externaldocument[algebraic-topology-]{algebraic-topology}
\externaldocument[differential-geometry-]{differential-geometry}
\externaldocument[algebraic-geometry-]{algebraic-geometry}
\externaldocument[symplectic-geometry-]{symplectic-geometry}
\externaldocument[complex-geometry-]{complex-geometry}
\externaldocument[hodge-theory-]{hodge-theory}
\maybeexternaldocument{deformations-}{deformations}
\externaldocument[vertex-algebras-]{vertex-algebras}
\externaldocument[springer-]{springer}
\externaldocument[infty-categories-]{infty-categories}
\externaldocument[seminars-differential-geometry-]{%
seminars-differential-geometry}
\externaldocument[seminars-complex-geometry-]{%
seminars-complex-geometry}
\externaldocument[seminars-algebraic-geometry-]{%
seminars-algebraic-geometry}
\externaldocument[seminars-others-]{%
seminars-others}
\externaldocument[%
geometric-prequantization-]{%
geometric-prequantization}
\externaldocument[surfaces-]{surfaces}
\externaldocument[birational-maps-commutative-algebra-]{birational-maps-commutative-algebra}
\externaldocument[cremona-transformations-]{cremona-transformations}
\externaldocument[derived-categories-of-sheaves-]{derived-categories-of-sheaves}
\externaldocument[geometric-invariant-theory-]{geometric-invariant-theory}
\externaldocument[geometric-stability-conditions-and-group-actions-]{geometric-stability-conditions-and-group-actions}
\externaldocument[moduli-spaces-of-sheaves-]{moduli-spaces-of-sheaves}
\externaldocument[bridgeland-stability-general-theory-]{bridgeland-stability-general-theory}
\externaldocument[hilbert-schemes-]{hilbert-schemes}
\externaldocument[noncommutative-abelian-surfaces-and-kummer-type-hyperkahler-varieties-]{noncommutative-abelian-surfaces-and-kummer-type-hyperkahler-varieties}
\externaldocument[mumford-tate-groups-in-hodge-theory-]{mumford-tate-groups-in-hodge-theory}
\externaldocument[hodge-birational-atoms-2026-]{hodge-birational-atoms-2026}
\externaldocument[xxii-egd-2026-]{%
xxii-egd-2026}
\externaldocument[pde-]{pde}
\externaldocument[probability-]{probability}

\externaldocument[linguistics-]{linguistics}
\externaldocument[genealogia-familiar-]{%
genealogia-familiar}

\externaldocument[%
16th-alga-meeting-2026-]{%
16th-alga-meeting-2026}

\externaldocument[%
iv-jornadas-lefschetz-2026-]{%
iv-jornadas-lefschetz-2026}

\externaldocument[reading-the-masters-]{%
reading-the-masters}

\externaldocument[real-analysis-]{%
real-analysis}

\externaldocument[optimal-transport-]{%
optimal-transport}

\externaldocument[history-]{history}

\externaldocument[seiberg-witten-]{seiberg-witten}

% Theorem environments.
%
\theoremstyle{plain}
\newtheorem{theorem}[subsection]{Theorem}
\newtheorem{proposition}[subsection]{Proposition}
\newtheorem{lemma}[subsection]{Lemma}

\theoremstyle{definition}
\newtheorem{definition}[subsection]{Definition}
\newtheorem{example}[subsection]{Example}
\newtheorem{exercise}[subsection]{Exercise}
\newtheorem{situation}[subsection]{Situation}

\theoremstyle{remark}
\newtheorem{remark}[subsection]{Remark}
\newtheorem{remarks}[subsection]{Remarks}

\numberwithin{equation}{subsection}

% Macros
%
\def\lim{\mathop{\mathrm{lim}}\nolimits}
\def\colim{\mathop{\mathrm{colim}}\nolimits}
\def\Spec{\mathop{\mathrm{Spec}}}
\def\Hom{\mathop{\mathrm{Hom}}\nolimits}
\def\Ext{\mathop{\mathrm{Ext}}\nolimits}
\def\SheafHom{\mathop{\mathcal{H}\!\mathit{om}}\nolimits}
\def\SheafExt{\mathop{\mathcal{E}\!\mathit{xt}}\nolimits}
\def\Sch{\mathit{Sch}}
\def\Mor{\mathop{\mathrm{Mor}}\nolimits}
\def\Ob{\mathop{\mathrm{Ob}}\nolimits}
\def\Sh{\mathop{\mathit{Sh}}\nolimits}
\def\NL{\mathop{N\!L}\nolimits}
\def\CH{\mathop{\mathrm{CH}}\nolimits}
\def\proetale{{pro\text{-}\acute{e}tale}}
\def\etale{{\acute{e}tale}}
\def\QCoh{\mathit{QCoh}}
\def\Ker{\mathop{\mathrm{Ker}}}
\def\Im{\mathop{\mathrm{Im}}}
\def\Coker{\mathop{\mathrm{Coker}}}
\def\Coim{\mathop{\mathrm{Coim}}}

% Boxtimes
%
\DeclareMathSymbol{\boxtimes}{\mathbin}{AMSa}{"02}

%
% Macros for moduli stacks/spaces
%
\def\QCohstack{\mathcal{QC}\!\mathit{oh}}
\def\Cohstack{\mathcal{C}\!\mathit{oh}}
\def\Spacesstack{\mathcal{S}\!\mathit{paces}}
\def\Quotfunctor{\mathrm{Quot}}
\def\Hilbfunctor{\mathrm{Hilb}}
\def\Curvesstack{\mathcal{C}\!\mathit{urves}}
\def\Polarizedstack{\mathcal{P}\!\mathit{olarized}}
\def\Complexesstack{\mathcal{C}\!\mathit{omplexes}}
% \Pic is the operator that assigns to X its picard group, usage \Pic(X)
% \Picardstack_{X/B} denotes the Picard stack of X over B
% \Picardfunctor_{X/B} denotes the Picard functor of X over B
\def\Pic{\mathop{\mathrm{Pic}}\nolimits}
\def\Picardstack{\mathcal{P}\!\mathit{ic}}
\def\Picardfunctor{\mathrm{Pic}}
\def\Deformationcategory{\mathcal{D}\!\mathit{ef}}

\begin{document}
\emergencystretch=1em
\title{VBAC}
\maketitle
\tableofcontents

VBAC webinar: Atoms and applications.
October 23, 2026, 13:00--15:00 UTC
(10:00--12:00 in Rio de Janeiro).
Organized by the VBAC Scientific Committee;
introduced and chaired by Ugo Bruzzo.
Webinar archive:
\url{http://vbac.wikidot.com/}.

\medskip\noindent
\textbf{Sources and editorial conventions.}
The abstracts are reproduced verbatim from
Lupita Hern\'andez's October 21 announcement.
The notes transcribe the supplied slide
screenshots as closely as possible;
repeated animation frames are consolidated.
Slide headings are retained as bold headings.
Formula layout and TeX typography are
adapted.
Statements below are recorded as presented
by the speakers, not independently verified
proofs. Oral additions are attributed,
timestamped paraphrases of the supplied Zoom
transcript; obvious speech-recognition errors
in mathematical names have been corrected.
The transcript labels Benedetti's speech
``julieth''; the introduction and continuity
of the presentation identify the speaker.
AI explanations are explicitly marked.

\section{Tony Pantev: Hodge atoms
and applications}
\label{section-vbac-pantev}
University of Pennsylvania.
October 23, 2026, 13:00--13:45 UTC
(10:00--10:45 Rio); discussion until 10:54.

\medskip\noindent
\textbf{Abstract (verbatim).}
I will explain how a natural amalgam of
classical Hodge theory with the nc Hodge
structures arising from Gromov-Witten theory
gives rise to new additive invariants of
smooth projective varieties called Hodge
atoms. Combined with Iritani's blow-up
formula, Hodge atoms provide obstructions
to birational equivalence. I will discuss
applications to classical rationality
problems. This is joint work with
L. Katzarkov, M. Kontsevich, and T. Y. Yu.

\medskip\noindent
\textbf{AI upshot.}
Combine classical Hodge theory with quantum
cohomology, split the resulting structure
into spectral pieces, and track those pieces
through blow-ups. A piece that cannot come
from projective space or a blow-up center
can obstruct rationality.

\medskip\noindent
\textbf{ncHodge structures.}
Screenshot 10:07:39.
A \emph{pure rational nc-Hodge structure}
is a triple $(H,\nabla,\mathcal E_B)$, where:
\begin{itemize}
\item $H\longrightarrow\mathbb C$ is a
$\mathbb Z/2$-graded algebraic
$\mathbb C$-vector bundle;
\item
$$
\mathcal E_B\subset
H|_{\mathbb C\setminus\{0\}}
$$
is a $\mathbb Q$-local system generating $H$;
\item $\nabla$ is a meromorphic connection
on $H$ with pole of order at most $2$ at
$0\in\mathbb C$.
\end{itemize}
Furthermore, $\mathcal E_B$ is
$\nabla$-flat, compatible with Stokes data,
and the vector bundle $\widehat H$ on
$\mathbb P^1$ obtained by gluing $H$ to its
complex conjugate is holomorphically trivial.

\medskip\noindent
\textbf{Oral explanation (10:05--10:14).}
Pantev explains that compatibility with
Stokes data means that the steps of the
Stokes filtration are defined over the
rational structure. The fiber at $1$ is the
underlying vector space; the connection
plays the role of a more complicated Hodge
filtration. Holomorphic triviality after
gluing expresses the analogue of the
opposedness of the Hodge filtration and its
complex conjugate. In response to the
question about ordinary Hodge structures,
he explains that periodizing by Tate twists
folds the integer grading into parity; the
resulting functor into nc-Hodge structures
is fully faithful. Ordinary Hodge theory
is thus retained after this periodization.

\medskip\noindent
\textbf{Remarks.}
Screenshot 10:11:15.
\begin{enumerate}
\item This is the structure arising
(or expected) on the cohomology of a
smooth and compact nc-space.
\item $\mathbb Z/2$-graded versions of the
classical pure Hodge structures from
K\"ahler geometry are incorporated here
as nc-Hodge structures with regular
singularities.
\item Families of nc-spaces give rise to
variations of nc-Hodge structures:
extension of the connection $\nabla$ to a
connection with regular singularities in
the direction of the parameter space of
the family.
\end{enumerate}

\medskip\noindent
\textbf{Quantum connection.}
Screenshot 10:13:41.
For a compact symplectic $(X,\omega)$ of
dimension $2d$ there is an associated
one-parameter variation of nc-Hodge
structures $(H,\nabla,\mathcal E_B)$ where
the Hodge filtration part $(H,\nabla)$ is
 given by
$$
H=H^\bullet(X,\mathbb C)[[q]][u,u^{-1}].
$$
The quantum connection is
$$
\begin{aligned}
\nabla &:H\longrightarrow H\otimes\Omega^1,\\
\nabla_{\partial_u}
&=\partial_u-u^{-2}(c_1(X)*_q(-))
+u^{-1}\operatorname{Gr}.
\end{aligned}
$$
Here $\operatorname{Gr}$ acts on
$H^k(X,\mathbb C)$ by
$$
\operatorname{Gr}=\frac{k-d}{2}
\operatorname{id},
$$
and the quantum product is
$$
*_q:H^\bullet(X,\mathbb C)^{\otimes2}
\longrightarrow H^\bullet(X,\mathbb C)
\otimes\mathbb C_q.
$$

\medskip\noindent
\textbf{Oral explanation (10:15--10:19).}
The variable $u$ encodes the nc-Hodge
filtration; $q$ is the variation parameter
scaling the symplectic form. In response to
a question about Gromov--Witten invariants,
Pantev locates them in $*_q$: its constant
term is the cup product, while higher terms
are Gromov--Witten corrections. This slide
uses the very small quantum product as an
example; the birational construction needs
the full quantum deformation.

\medskip\noindent
\textbf{$F$-bundles.}
Screenshot 10:19:09.
HMS lesson: nc-Hodge structures are a
powerful tool for linearizing various
geometries.

Downside: nc-Hodge structures are difficult
to compute and classify.

Goal: find a tractable and a computationally
viable shadow of a variation of nc-Hodge
structures.

This leads to the notion
[Katzarkov--Kontsevich--Pantev--Yu, 2020]
of an $F$-bundle, which turns out to be
more user friendly.

\medskip\noindent
\textbf{Oral motivation (10:19--10:28).}
The categories on the two sides of mirror
symmetry are difficult geometric objects.
Passing to their cohomology linearizes them,
but the nc-Hodge structure is still hard to
compute. An $F$-bundle keeps a more tractable
part of this structure, discarding the
rational structure and opposedness condition.

Analytic continuation is needed to compare
the relevant objects at different points.
Complex convergence of Gromov--Witten series
is expected in this setting but is known
only in limited cases. Pantev therefore
works with rigid analytic spaces over a
non-Archimedean field: the required
convergence and continuation are available
there. The choice of field addresses a real
technical obstruction, not just notation.

\medskip\noindent
\textbf{$F$-bundles (ii).}
Screenshot 10:23:56.
Fix a non-Archimedean normed field
$\mathbb k$ of characteristic zero.
Write $\mathbb D$ for the germ at $0$ of a
$\mathbb k$-analytic disk with variable $u$.
Let $o\in\mathbb D$ be the closed point.

An \emph{$F$-bundle} is the data
$(\mathcal H,\nabla)/\mathcal B$, where
$\mathcal B$ is a $\mathbb k$-analytic
supermanifold, $\mathcal H$ is a super
vector bundle over
$\mathcal B\times\mathbb D$, and $\nabla$
is a connection with poles at most at
$u=0$, satisfying:
\begin{itemize}
\item $\nabla^2=0$;
\item $\nabla_{u^2\partial_u}$ is regular
(no poles);
\item $\nabla_{u\xi}$ is regular for every
$\xi\in T\mathcal B$.
\end{itemize}

\medskip\noindent
\textbf{$F$-bundles (iii): residual
structures.}
Screenshot 10:29:43.
An $F$-bundle induces:
\begin{itemize}
\item A vector bundle $\mathcal H_o$ on
$\mathcal B$.
\item A relative flat $u$-connection
$(\mathcal H,\delta)$ on
$\mathcal B\times\mathbb D\to\mathbb D$.
\item A regular Higgs field
$$
\theta:\mathcal H_o\longrightarrow
\mathcal H_o\otimes\Omega^1_{\mathcal B},
\qquad \theta\wedge\theta=0.
$$
\item A regular endomorphism
$$
\kappa:\mathcal H_o\longrightarrow
\mathcal H_o
$$
which supercommutes with $\theta$.
\end{itemize}
Here
$$
\mathcal H_o=
\mathcal H|_{\mathcal B\times\{o\}}.
$$
The relative operator
$$
\delta:\mathcal H\longrightarrow
\mathcal H\otimes\Omega^1_{\mathcal B}\{u\}
$$
is the projection of $u\nabla$;
$\theta$ is its restriction to $u=0$;
$\kappa$ is the restriction of
$\nabla_{u^2\partial_u}$ to $u=0$.

\medskip\noindent
\textbf{Oral explanation (10:28--10:32).}
Multiplying the connection by $u$ changes
its Leibniz rule: its principal symbol is
$u$ times the identity. At $u=0$ the
derivative part disappears, leaving an
$\mathcal O$-linear Higgs field. The leading
coefficient in the $u$ direction similarly
becomes the endomorphism $\kappa$.
Flatness supplies their compatibility.

\medskip\noindent
\textbf{$F$-bundles (iv): maximality.}
Screenshot 10:33:57.
The associated Higgs field, viewed as a map
$$
\theta\mathbin{\lrcorner}(-):
T_{\mathcal B}\longrightarrow
\operatorname{End}(\mathcal H_o),
$$
realizes the vector fields on $\mathcal B$
as a commuting family of
$\mathcal O_{\mathcal B}$-linear operators.
We say the $F$-bundle is \emph{maximal at
$b\in\mathcal B$} if there exists an even
vector $h\in(\mathcal H_o)_b$ such that
$$
\begin{aligned}
T_{\mathcal B,b}&\longrightarrow
(\mathcal H_o)_b,\\
\xi&\longmapsto
(\theta_b\mathbin{\lrcorner}\xi)(h)
\end{aligned}
$$
is an isomorphism. It is \emph{maximal} if
it is maximal at every point.

\textbf{Oral explanation (10:33).}
Maximality means the commuting operators
have a common cyclic vector; its existence,
rather than a choice of it, is required.

\medskip\noindent
\textbf{$F$-bundles (v): properties.}
Screenshot 10:35:10.
Suppose the $F$-bundle is maximal. Then:
\begin{enumerate}
\item $T_{\mathcal B}$ inherits the structure
of a unital, commutative, associative algebra
naturally embedded in
$\operatorname{End}(\mathcal H_o)$.
Hence $\mathcal B$ is a $\mathbb k$-analytic
$F$-manifold in the sense of
Hertling--Manin.
\item $\mathcal B$ has an Euler vector field,
the unique vector field
$$
\operatorname{Eu}\in
\Gamma(\mathcal B,T_{\mathcal B})
$$
such that
$$
\theta\mathbin{\lrcorner}
\operatorname{Eu}=\kappa.
$$
Its powers with respect to the product on
$T_{\mathcal B}$ generate an action of the
Witt Lie algebra on $\mathcal B$.
\end{enumerate}

\medskip\noindent
\textbf{$F$-bundles (vi): decomposition.}
Screenshot 10:37:07.
Let $b\in\mathcal B$ be a rigid point.
The decomposition
$$
\mathcal H_{(b,o)}=
\bigoplus_\lambda\mathcal H^\lambda_{b,o}
$$
into generalized eigenspaces for the
distinct eigenvalues of $\kappa_b$ induces
a canonical decomposition of the germ of
the maximal $F$-bundle into an external
product of maximal $F$-bundles
$$
\mathop{\boxplus}_{
\lambda\in\operatorname{spec}(\kappa_b)}
(\mathcal H^\lambda,\nabla^\lambda)
/\mathcal B^\lambda.
$$
Slide attribution:
Katzarkov--Kontsevich--Pantev--Yu (2021),
formal case; Hinault--Yu--Zhang--Zhang
(2024), $\mathbb k$-analytic case.

\textbf{Oral clarification (10:35--10:37;
10:45--10:46).}
This is stronger than decomposing one
vector space: the base decomposes locally
as a product, the bundle as a sum of
pullbacks from those factors, and the
connection as the corresponding block
sum. The pieces retain their irregular
connections; decomposition does not make
the connections regular singular.

\medskip\noindent
\textbf{$F$-bundles (vi): blow-up theorem.}
Screenshots 10:38:41 and 10:40:28.
The main example is the A-model $F$-bundle
$$
({}^X\mathcal H,{}^X\nabla)/{}^X\mathcal B
$$
corresponding to the universal family of
Fukaya categories of a compact symplectic
manifold $X$.

Suppose $X$ is smooth projective over
$\mathbb C$, $Y\subset X$ is smooth of
codimension $m\geq2$, and
$\widehat X=\operatorname{Bl}_Y(X)$.
There are special limiting points in
$$
{}^{\widehat X}\mathcal B
\qquad
{}^X\mathcal B\times
({}^Y\mathcal B)^{\times(m-1)}
$$
such that the germs of $F$-bundles satisfy
$$
\begin{aligned}
({}^{\widehat X}\mathcal H,
{}^{\widehat X}\nabla)
\cong{}&({}^X\mathcal H,{}^X\nabla)\\
&\boxplus
({}^Y\mathcal H,{}^Y\nabla)
^{\boxplus(m-1)}.
\end{aligned}
$$
The germs are taken at their respective
limiting points. Slide attribution:
Iritani (2023), KKPY (2024).

\textbf{Oral explanation (10:37--10:41).}
The A-model connection is the full quantum
connection, allowing cohomology parameters
in all degrees. The formula parallels the
ordinary cohomology formula for a blow-up,
but retains the quantum data. The two
limiting points are different: comparing
them is exactly where analytic continuation
and the rigid analytic formulation enter.

\medskip\noindent
\textbf{$F$-bundles (vii).}
Screenshot 10:42:29.
The blow-up formula is compatible with
spectral decompositions on both sides.
Thus the pieces in a spectral decomposition
of the A-model $F$-bundle can serve as
blow-up invariants. To formalize this we
must keep track of these pieces together
with their natural symmetries.

\medskip\noindent
\textbf{Mumford--Tate groups.}
Screenshots 10:42:59 and 10:44:12.
Consider $\operatorname{MT}$, the
Mumford--Tate group: the Galois group of the
neutral Tannakian category of polarizable
pure rational Hodge structures.
The homomorphisms
$$
\iota:\mathbb G_m\longrightarrow
\operatorname{MT},\qquad
\pi:\operatorname{MT}\longrightarrow
\mathbb G_m
$$
correspond to the Hodge grading and the
action on $H^2(\mathbb P^1)$ respectively.
Set
$$
\operatorname{Hod}=\ker\pi,\qquad
\epsilon_{\operatorname{Hod}}=\iota(-1).
$$
The marked element has order two.
The group $\operatorname{Hod}$ is
proreductive and is the Galois group of the
Tannakian category of pure polarizable
$\mathbb Z/2$-graded rational Hodge
structures. For every smooth complex
projective $X$, its marked element acts on
$H^i_{\mathrm{dR}}(X,\mathbb C)$ by $(-1)^i$.

\medskip\noindent
\textbf{Spectral covers.}
Screenshots 10:44:44 and 10:45:41.
Let $\mathbb k\supset\mathbb Q$ be an
algebraically closed non-Archimedean
normed field whose norm on $\mathbb Q$
is trivial, for example the field denoted
on the slide by
$$
\mathbb k=\overline{\mathbb Q}
((y^{\mathbb Q})).
$$
Let $(\mathcal H,\nabla)/\mathcal B_X$ be
the $\mathbb k$-analytic A-model maximal
$F$-bundle of $X$. The base is an
admissible open in
$H^\bullet(X,\mathbb k)^{\mathrm{an}}$.
The group
$\operatorname{Hod}(\overline{\mathbb Q})$
acts by analytic automorphisms.
Its fixed points
$\mathcal B_X^{\operatorname{Hod}}$ form a
purely even $\mathbb k$-analytic manifold,
an admissible open in the analytification
of the rational Hodge classes tensored with
$\mathbb k$.

The residual operator
$$
\kappa=\operatorname{Eu}*(-)
$$
on $\mathcal H_o$ restricted to this locus
defines a finite spectral cover
$$
\Sigma_X\longrightarrow
\mathcal B_X^{\operatorname{Hod}}.
$$
It is the zero locus of
$\det(\lambda\operatorname{id}-\kappa)$ in
$$
\mathcal B_X^{\operatorname{Hod}}
\times\mathbb A^{1,\mathrm{an}}_{\mathbb k}.
$$
Consider its reduction
$\Sigma_{X,\mathrm{red}}$, the complement
$\mathcal U_X$ of the branch divisor in
$\mathcal B_X^{\operatorname{Hod}}$, and
$$
\widetilde{\mathcal U}_X=
\mathcal U_X\times_{
\mathcal B_X^{\operatorname{Hod}}}
\Sigma_{X,\mathrm{red}}.
$$
Then
$\widetilde{\mathcal U}_X\to\mathcal U_X$
is an unramified $\mathbb k$-analytic cover.

\textbf{Oral conclusion (10:44--10:45).}
Pantev identifies the connected components
of the spectral cover, with their associated
pieces, as the atoms. He mentions the
irrationality of very general cubic
fourfolds as an application, but does not
present the example because time runs out.
The blow-up formula adds pieces from the
center; it does not assert that the entire
collection is unchanged by every blow-up.

\medskip\noindent
\textbf{Discussion not on the slides
(10:46--10:54).}
\begin{itemize}
\item Asked about the scope, Pantev says
the construction applies to all smooth
projective varieties.
\item Asked about higher-genus
Gromov--Witten invariants, he distinguishes
them from the genus-zero data used here.
He describes higher-genus information in
terms of quantization of moduli spaces of
objects and the BCOV formalism. For
Calabi--Yau categories these moduli spaces
carry shifted symplectic structures;
more generally Poisson structures enter.
He stresses that identifying the specific
geometric quantization selected by the
higher-genus theory remains unclear.
This is his explanatory discussion, not a
proved extension of the atom construction.
\item Gabriel asks about a genus-one
analogy involving the Quillen metric and
analytic torsion. Pantev says one must
first understand how higher-genus theory
quantizes the moduli of objects.
\item Kai asks whether the decomposition's
$R$-matrix reconstructs higher-genus theory
of a blow-up from that of the base.
Pantev explicitly answers that he does not
know; no affirmative claim is recorded.
\end{itemize}

\section{Vladimiro Benedetti: An atomic
criterion for irrationality of
hyperplane sections}
\label{section-vbac-benedetti}
Universit\'e C\^ote d'Azur, Nice.
October 23, 2026, 14:00--14:45 UTC
(11:00--11:45 Rio); discussion until 11:55.

\medskip\noindent
\textbf{Abstract (verbatim).}
In this talk, I will give some examples of
applications of the theory of Hodge atoms
in birational geometry. More precisely,
I will explain how one can get irrationality
of very general hyperplane sections of some
Fano fivefolds, such as Gushel-Mukai or
Kuchle fourfolds of type (c5). In order to
do so, I will recall the interest in
studying these varieties, and how the
monodromy invariance of quantum cohomology
allows to deduce an irrationality criterion
for fourfolds without having to explicitly
compute the quantum ring. This is a joint
work with J\'er\'emy Gu\'er\'e,
Laurent Manivel and Nicolas Perrin.

\medskip\noindent
\textbf{AI comment: comparison with Atoms.}
The earlier talk, ``Quantum cohomology and
irrationality of Gushel-Mukai fourfolds'',
is recorded in the chapter
\emph{Hodge theory, birational geometry
and atoms 2026},
\href{https://www.heap-project.org/tag/04YR}
{section on Gushel--Mukai fourfolds}.
Both talks use quantum spectral pieces,
Hodge structures and the blow-up formula
to obstruct rationality. The earlier talk
focused on computing the Gushel--Mukai
quantum ring and on associated K3 surfaces.
Today's talk replaces those explicit
computations by monodromy and a numerical
bound for surface atoms, giving a common
criterion for cubics, Gushel--Mukai and
K\"uchle type (c5) fourfolds. Thus the
strategy is closely related, but the
criterion and its scope are broader.

\medskip\noindent
\textbf{Results on cubics.}
Screenshots 11:02:04 and 11:02:30, slide 3.
Theory of atoms over $\mathbb Q$
[Katzarkov--Kontsevich--Pantev--Yu].

\begin{theorem}[KKPY, as presented]
\label{theorem-vbac-cubic}
The very general cubic fourfold is not
rational.
\end{theorem}

\begin{theorem}[Gu\'er\'e, as presented]
\label{theorem-vbac-guere}
If $X$ is a rational cubic fourfold, then
there exists a projective K3 surface $S$
and an isomorphism of rational Hodge
structures
$$
H^4(X,\mathbb Q)_{\mathrm{van}}
\cong H^2(S,\mathbb Q)(-1).
$$
\end{theorem}
The slide's key point is that birational
behaviour is determined by Hodge structure,
with the displayed decomposition
$$
H^4=H^{1,3}\oplus H^{2,2}\oplus H^{3,1}.
$$

\medskip\noindent
\textbf{Results on hyperplane sections.}
Screenshot 11:04:24, slide 4.

\begin{theorem}[Benedetti--Fay--Gu\'er\'e--
Manivel--Perrin, as presented]
\label{theorem-vbac-hyperplanes}
Let $Y\subset\mathbb P^n$ be a smooth
Fano fivefold. Let $X$ be a smooth Fano
hyperplane section, and suppose that
$$
\begin{gathered}
b_1(X)=b_3(X)=0,\\
p_g(X):=h^{3,1}(X)>h^{3,1}(Y),\\
b_4(X)_{\mathrm{van}}\geq10+12p_g(X).
\end{gathered}
$$
If $X$ is Hodge general, it is irrational.
\end{theorem}
Here $p_g(X)$ is the speaker's notation
for $h^{3,1}(X)$, not the usual geometric
genus $h^{4,0}(X)$ of a fourfold.
The slide credits Aideen Fay as well as
the collaborators named in the abstract.

Applies to:
\begin{itemize}
\item Cubic fourfolds in $\mathbb P^5$
[KKPY].
\item Gushel--Mukai fourfolds
$$
\begin{aligned}
X&=G(2,5)\cap H\cap Q,\\
X&\subset G(2,5)\cap Q
\end{aligned}
$$
[Benedetti--Manivel--Perrin].
\item K\"uchle fourfolds of type (c5).
\end{itemize}
No computation in quantum cohomology
necessary.

\medskip\noindent
\textbf{Hodge generality.}
Screenshots 11:07:08 and 11:08:01, slide 5.
For $\iota:X\hookrightarrow Y$, put
$$
\begin{aligned}
H^\bullet(X)_{\mathrm{amb}}
&=\iota^*H^\bullet(Y),\\
H^\bullet(X)_{\mathrm{van}}
&=\ker(\iota_*),\\
H^4(X)_{\mathrm{Hdg}}
&=H^{2,2}(X)\cap H^4(X,\mathbb Q).
\end{aligned}
$$
The section is \emph{Hodge general} if
$$
H^4(X)_{\mathrm{Hdg}}
=\iota^*H^4(Y)_{\mathrm{Hdg}}.
$$
The slide writes $j^*$ in this last formula;
here it is written $\iota^*$ consistently.

\begin{proposition}[As presented]
\label{proposition-vbac-hodge-general}
Suppose $X$ is a very general hyperplane
section and $h^{3,1}(Y)<h^{3,1}(X)$.
Then $X$ is Hodge general.
\end{proposition}
\begin{proof}
For $X$ very general, its rational Hodge
classes are stable under monodromy
[Voisin, as cited on the slide].
They cannot meet the vanishing cohomology
nontrivially: otherwise they contain it,
since its monodromy representation is
irreducible. But
$$
H^{3,1}(X)\cap H^4(X)_{\mathrm{van}}
\neq0.
$$
\end{proof}

\textbf{Oral explanation (11:07--11:09).}
A loop in the parameter space of smooth
hyperplane sections gives an automorphism
of cohomology by parallel transport.
The irreducible representation here is
specifically the vanishing cohomology, not
all of cohomology. The inequality on
$h^{3,1}$ guarantees a non-Hodge direction
in that representation, ruling out its
being entirely made of Hodge classes.

\medskip\noindent
\textbf{Period map of cubics, GM, (c5).}
Screenshots 11:10:33 and 11:11:15, slide 6.
For these fourfolds, the Hodge structure
splits into ambient and vanishing parts.
The slide's even-degree Hodge diamonds are
$$
\begin{array}{ccccc}
&&1&&\\
&0&1&0&\\
0&1&20+*&1&0\\
&0&1&0&\\
&&1&&
\end{array}
=
\begin{array}{ccccc}
&&1&&\\
&0&1&0&\\
0&0&*&0&0\\
&0&1&0&\\
&&1&&
\end{array}
\oplus
\begin{array}{ccccc}
&&0&&\\
&0&0&0&\\
0&1&20&1&0\\
&0&0&0&\\
&&0&&
\end{array}.
$$
Set
$$
\Lambda=(H^4(X,\mathbb Z)_{\mathrm{van}},
\cup).
$$
The period map sends
$$
[X]\longmapsto[H^{3,1}(X)]
\in\mathbb P(\Lambda\otimes_{\mathbb Z}
\mathbb C).
$$
The slide denotes this map by
$$
p:\mathcal H_{X\subset Y}
\longrightarrow\mathcal P.
$$
\begin{itemize}
\item Cubics: $p$ is birational [Voisin].
\item GM: fibers of $p$ are four-dimensional,
but GMs in each fiber are birational
[Debarre--Kuznetsov].
\item (c5): less is known.
\end{itemize}
These are the slide's moduli-level claims;
its notation for the parameter space is
schematic, rather than a specification of
markings and quotients.

\medskip\noindent
\textbf{Quantum cohomology I.}
Screenshot 11:13:24, slide 7.
Let $d\in\operatorname{NE}(X)
\subset H_2(X,\mathbb Z)$ be an effective
curve class. Let $\overline M_{0,n}(X,d)$
be the space of genus-zero stable maps of
degree $d$ with $n$ markings.
The evaluation maps are
$$
\begin{aligned}
\operatorname{ev}_i:
\overline M_{0,n}(X,d)&\longrightarrow X,\\
f&\longmapsto f(p_i).
\end{aligned}
$$
Expected dimension:
$$
\operatorname{edim}_{d,n}(X)
=\dim X+\int_d c_1(X)+n-3.
$$
For $\sigma_i\in H^\bullet(X)$, the
Gromov--Witten invariant is
$$
\left\langle\sigma_1,\ldots,\sigma_n
\right\rangle^X_{d,n}
=
\int_{[\overline M_{0,n}(X,d)]
^{\mathrm{virt}}}
\prod_{i=1}^n\operatorname{ev}_i^*(\sigma_i).
$$
GW invariants are deformation invariant.

\textbf{Oral explanation (11:13--11:15).}
The domain of a stable map may be a nodal
genus-zero curve with marked points.
The integral uses the virtual fundamental
class of the expected dimension; one cannot
in general replace it by the ordinary
fundamental class. Deformation invariance
will be used as monodromy invariance.

\medskip\noindent
\textbf{Quantum cohomology II.}
Screenshot 11:17:01, slide 8.
Let $(\Delta_i)$ be a graded basis of
$H^\bullet(X,\mathbb Q)$ and $(\Delta^i)$
the dual basis. Use dual coordinates $t_i$,
with
$$
\deg(t_i)=2-\deg(\Delta_i),\qquad
\tau=\sum_i t_i\Delta_i.
$$
Work in the Novikov ring, with
$$
\deg(q^d)=\int_d c_1(X).
$$
The quantum product on
$$
QH(X)=H^\bullet(X,\mathbb Q)
\otimes_{\mathbb Q}\mathbb Q[[q,(t_j)]]
$$
is
$$
\Delta_i\star_\tau\Delta_j
=
\sum_{d\geq0,n\geq0}\frac{q^d}{n!}
\sum_k
\left\langle\Delta_i,\Delta_j,\Delta_k,
\tau,\ldots,\tau\right\rangle^X_{d,n+3}
\Delta^k.
$$
The Euler vector field is
$$
\operatorname{Eu}_\tau=c_1(X)
-\sum_i\frac{\deg(t_i)t_i}{2}\Delta_i.
$$
The main character is
$$
\kappa_\tau=
\operatorname{Eu}_\tau\star_\tau(-).
$$
The grading and sign convention above
follow the slide.

\medskip\noindent
\textbf{Euler vector field invariants.}
Screenshot 11:20:17, slide 9.
Fix coordinates of a point
$\tau\in H^\bullet(X,\mathbb Q)
_{\mathrm{Hdg}}$
so that $\kappa_\tau$ has maximal number
of eigenvalues. Its spectrum is the set
$\operatorname{Sp}(\kappa_\tau)$ of these
eigenvalues. For each $\alpha$ in it, put
$$
E_{\tau,\alpha}=
\ker(\kappa_\tau-\alpha\operatorname{id})^N,
\qquad N\gg0.
$$
The slide points to evaluation functions
$$
\operatorname{ev}_\tau:
\mathbb Q[[q,(t_j)]]\longrightarrow
F[t^{\pm1}]
$$
to a Levi--Civita field $F$.

\textbf{Oral qualification (11:18--11:20).}
Benedetti explicitly calls the language
``fix a Hodge class'' a simplification.
Because convergence is an issue, the proper
construction uses a suitable evaluation
function and takes the spectrum of the
evaluated operator. Maximality is with
respect to the allowed evaluations.

\medskip\noindent
\textbf{Iritani's blowup formula I.}
Screenshot 11:22:19, slide 10.
Let
$$
\pi:\widetilde X=\operatorname{Bl}_Z(X)
\longrightarrow X,
\qquad c=\operatorname{codim}_X(Z).
$$
There is an isomorphism of Hodge structures
$$
H^\bullet(\widetilde X,\mathbb Q)
\cong H^\bullet(X,\mathbb Q)
\oplus\bigoplus_{k=1}^{c-1}
H^{\bullet-2k}(Z,\mathbb Q)(-k).
$$
Let $q$, $\widetilde q$, and $q_Z$ be the
Novikov variables on $X$, $\widetilde X$,
and $Z$. They are embedded into a common
Laurent-series ring
$$
K[q_E^{\pm1/r}][[q]].
$$

\textbf{Oral explanation (11:21--11:23).}
Here $K$ extends $\mathbb Q$, and $q_E$
corresponds to a curve in the exceptional
divisor. The important point is that
negative powers of $q_E$ are needed;
the talk leaves the integer $r$ unspecified.

\medskip\noindent
\textbf{Iritani's blowup formula II.}
Screenshot 11:24:03, slide 11.
Write $(\widetilde t_j)$ for the dual
variables on $\widetilde X$.
There exist an extension
$\mathbb Q\subset K$ and power series
$$
\begin{aligned}
f&\in H^\bullet(X,\mathbb Q)
\otimes_{\mathbb Q}
K[q_E^{\pm1/r}][[q,(\widetilde t_j)]],\\
f_k&\in H^\bullet(Z,\mathbb Q)
\otimes_{\mathbb Q}
K[q_E^{\pm1/r}][[q,(\widetilde t_j)]],
\end{aligned}
$$
for $1\leq k\leq c-1$, and an isomorphism
over this ring
$$
\begin{aligned}
\Psi:(QH^\bullet(\widetilde X),
\star_{\widetilde\tau})
\longrightarrow{}&
(QH^\bullet(X),\star_{f(\widetilde\tau)})\\
&\oplus\bigoplus_{k=1}^{c-1}
(QH^{\bullet-2k}(Z),
\star_{f_k(\widetilde\tau)}).
\end{aligned}
$$
It preserves Hodge structures over $K$
and Hochschild degree, and
$$
\Psi\kappa_{\widetilde\tau}\Psi^{-1}
=
\kappa^X_{f(\widetilde\tau)}
\oplus\bigoplus_{k=1}^{c-1}
\kappa^Z_{f_k(\widetilde\tau)}.
$$
The slide uses a sum sign for this last
block operator; direct-sum notation makes
its action on the preceding decomposition
explicit.

\textbf{Oral explanation (11:23--11:26).}
Think of $f$ and $f_k$ as changes of
coordinates relating the quantum parameters
on the blow-up, the base and the center.
They preserve the Hodge locus.
At first, a generalized eigenspace could
contain contributions from both base and
center if their eigenvalues coincide.
The next statement prevents this mixing
at the relevant maximal-spectrum points.

\medskip\noindent
\textbf{Iritani's blowup formula: remarks.}
Screenshots 11:28:25 and 11:30:39, slide 12.
The Novikov ring depends on the blow-up
(the inclusion of $q_Z$ in $q$ and the
need to invert $q_E$). This restricts
possible evaluations.

Proposition [Gu\'er\'e, KKPY, as presented]:
fix coordinates so that
$\kappa_{\widetilde\tau}^{\widetilde X}$
has maximal number of eigenvalues. Then
\begin{itemize}
\item $\kappa^X_{f(\widetilde\tau)}$ has
maximal number of eigenvalues;
\item the spectra of this operator and of
$$
\kappa^Z_{f_1(\widetilde\tau)},\ldots,
\kappa^Z_{f_{c-1}(\widetilde\tau)}
$$
are pairwise disjoint.
\end{itemize}
For the center, maximality is only with
respect to the blow-up, not necessarily
among all parameters for $Z$.

\textbf{Oral intuition and final question
(11:29--11:31; 11:52--11:55).}
A collision between a base eigenvalue and
a center eigenvalue, or between different
center copies, can be separated by shifting
a center parameter in the unit direction.
By the fundamental-class axiom this shifts
the corresponding operator by a scalar
multiple of the identity. A collision
would therefore contradict maximality of
the number of distinct eigenvalues.
Here ``fundamental class'' refers to the
unit $\Delta_0=1$ in cohomology, not the
point class. With the displayed convention
for $\operatorname{Eu}$, the scalar shift
has a minus sign; the separation argument
uses only that it is a scalar shift.

An eigenvalue at a maximal-spectrum point
belongs to a connected component of Tony's
spectral cover. The generalized eigenspace
is the linear data associated to that atom.

\medskip\noindent
\textbf{Oral ingredient missing from the
screenshots (11:31--11:33).}
Benedetti states that if $K_X$ is nef,
then, at the permitted Hodge parameters,
$\kappa_\tau$ differs from a scalar operator
by a nilpotent operator. Consequently all
cohomology belongs to a single generalized
eigenspace. He says this follows from the
degrees of the Novikov variables.
The transcript's sign for the scalar is
not reliable, so only the sign-independent
statement is recorded here.

\medskip\noindent
\textbf{Atoms of surfaces: key properties.}
Screenshot 11:35:50, slide 14.
Let $S$ be a projective surface,
$\tau\in H^\bullet(S)_{\mathrm{Hdg}}$, and
$\alpha\in\operatorname{Sp}(\kappa_\tau)$.
The notation for Hochschild degree is
$$
H^{(i)}(S)=\bigoplus_{p-q=i}H^{p,q}(S).
$$
Proposition [BFGMP, as presented]: if
$$
E_{\tau,\alpha}\cap H^{(2)}(S)\neq0,
\qquad
E_{\tau,\alpha}\cap H^{(1)}(S)=0,
$$
then
$$
\operatorname{codim}_{E_{\tau,\alpha}}
(E_{\tau,\alpha}\cap
H^\bullet(S)_{\mathrm{Hdg}})
\leq9+12p_g(S).
$$
Here $p_g(S)=h^{2,0}(S)$.

Proof on the slide: consider a minimal
model $S\to\Sigma$, with $S$ obtained by
blowing up finitely many points $Z$.
The blow-up formula conjugates the operator
to the sum of the operators for $\Sigma$
and the centers. Since
$H^{(2)}(S)\neq0$, the minimal model has
nef canonical class, so its cohomology
belongs to one generalized eigenspace.
The pieces for $Z$ are generated by Hodge
classes. The hypotheses imply
$c_1(\Sigma)^2\geq0$, and Riemann--Roch
implies
$$
h^{1,1}(\Sigma)\leq10+10p_g(\Sigma).
$$

\textbf{Oral explanation (11:34--11:36).}
Point blow-ups contribute only Hodge
classes and therefore do not increase the
codimension of the Hodge classes in the
relevant piece. The issue reduces to the
minimal surface and a numerical bound on
its cohomology.

\medskip\noindent
\textbf{Computation of GW invariants with
ambient insertions.}
Screenshot 11:39:37, slide 15.
For complete intersections $X\subset Y$,
techniques include:
\begin{itemize}
\item Homogeneous $Y$: GW invariants as
counts of curves.
\item Quantum Lefschetz by Coates and
Givental.
\item GKM varieties with torus action:
localization formulas. The slide cites
Holmes--Muratore and
Alexeev--Cavenaghi--Galindo--Georgiev--
Katzarkov--Martins.
\end{itemize}
These give invariants with insertions from
$H^\bullet(X)_{\mathrm{amb}}$, not the
vanishing classes.

\textbf{Oral addition (11:37--11:39).}
Benedetti mentions software from the last
two groups for GKM computations. The
limitation is decisive: computing ambient
insertions alone does not recover the
whole quantum ring. Monodromy provides the
needed information about the vanishing part.

\medskip\noindent
\textbf{Monodromy action's consequences.}
Screenshot 11:40:35, slide 16.
Let $\iota:X\hookrightarrow Y$ be a
hyperplane section. GW invariants are
monodromy invariant, and the vanishing
cohomology is an irreducible monodromy
representation.
If all $\sigma_i$ are ambient and $\gamma$
is vanishing, then
$$
\left\langle\sigma_1,\ldots,\sigma_k,
\gamma\right\rangle^X_{d,k+1}=0.
$$
For ambient $\alpha$ and $\tau$ this implies
$$
\begin{aligned}
\alpha\star_\tau
QH^\bullet(X)_{\mathrm{amb}}
&\subset QH^\bullet(X)_{\mathrm{amb}},\\
\alpha\star_\tau
QH^\bullet(X)_{\mathrm{van}}
&\subset QH^\bullet(X)_{\mathrm{van}}.
\end{aligned}
$$
Moreover, the action on the vanishing part
is scalar; the same holds for $\kappa_\tau$:
$$
\kappa_\tau=
\begin{pmatrix}
\kappa_\tau|_{\mathrm{amb}}&0\\
0&\lambda\operatorname{id}
\end{pmatrix}.
$$

\textbf{Oral emphasis (11:39--11:40).}
The entire vanishing cohomology lies in
one generalized eigenspace. For the
irrationality argument it is enough to
know this; one need not compute the
ambient block or the scalar $\lambda$.

\medskip\noindent
\textbf{The case of (c5).}
Screenshot 11:42:13, slide 17.
We have
$$
X\hookrightarrow Y\subset G(3,7).
$$
The fourfold $X$ is the zero locus of a
section of
$$
\mathcal O(1)\oplus\Lambda^2\mathcal U^\vee
\oplus\mathcal Q^\vee(1),
$$
and the fivefold $Y$ is the zero locus of a
section of
$$
\Lambda^2\mathcal U^\vee
\oplus\mathcal Q^\vee(1).
$$
Problems listed on the slide:
\begin{itemize}
\item $Y$ deforms, so the period map is
more complicated.
\item No rational smooth examples were
known to the speaker.
\item GW invariants with insertions from
$\iota^*H^\bullet(G(3,7))$ are accessible,
but these do not include all classes
coming from $Y$.
\item The vanishing block is scalar, but
$\lambda$ is unknown (index one).
\end{itemize}

\medskip\noindent
\textbf{Proof of the main theorem.}
Screenshot 11:44:03, slide 18.
Take a Hodge parameter $\tau$ for which
$\kappa_\tau$ has maximal number of
eigenvalues. Monodromy gives
$$
QH^\bullet(X)_{\mathrm{van}}
\subset E_{\tau,\lambda}.
$$
Consequently
$$
\begin{gathered}
p_g(X)\geq
\dim(E_{\tau,\lambda}\cap H^{(2)}(X))
\geq h^{3,1}(X)-h^{3,1}(Y)>0,\\
E_{\tau,\lambda}\cap H^{(1)}(X)=0.
\end{gathered}
$$
By Hodge generality,
$$
\operatorname{codim}_{E_{\tau,\lambda}}
(E_{\tau,\lambda}\cap
H^\bullet(X)_{\mathrm{Hdg}})
\geq10+12p_g(X).
$$
The surface bound rules out this piece
coming from a surface; points and curves
cannot supply its nonzero Hochschild
 degree-two part either.

\textbf{Oral explanation (11:43--11:45).}
If the piece came from a surface $S$, its
Hochschild degree-two dimension would bound
$p_g(S)$ by $p_g(X)$, whereas its Hodge
codimension would exceed the surface bound.
Thus the numbers are incompatible before
one knows the detailed quantum ring.

\medskip\noindent
\textbf{Weak factorization.}
Screenshot 11:47:20, slide 19.
Assume $X$ is birational to $\mathbb P^4$.
Weak factorization gives a chain of
blow-ups and blow-downs through smooth
varieties, starting at $X_0=X$ and ending
at $X_q=\mathbb P^4$.
\begin{itemize}
\item Start from $X_0$ and move towards
$X_q$, applying the blow-up formula.
\item Inductively construct Hodge parameters
$\tau_i$ with maximal spectrum and pieces
$E_{\tau_i,\lambda_i}$ having the same
invariants as $E_{\tau,\lambda}$.
\item This is possible because these
invariants cannot occur for points, curves
or surfaces, and base pieces do not merge
with pieces from blow-up centers.
\item Contradiction at $\mathbb P^4$.
\end{itemize}

\textbf{Oral intuition (11:45--11:47).}
All nontrivial smooth blow-up centers in a
fourfold have dimension at most two.
The distinguished piece cannot disappear
into any such center, so it survives along
the whole chain. But projective space has
no piece with its nonzero Hochschild
degree-two contribution.

\medskip\noindent
\textbf{Discussion not on the slides
(11:47--11:51).}
\begin{itemize}
\item Asked where the generalized eigenspace
comes from, Benedetti reiterates that one
first chooses a Hodge parameter, obtains
$\kappa_\tau$, then takes its generalized
eigenspaces. The maximal-spectrum condition
is the additional choice needed here.
\item Other irrationality criteria use the
same weak-factorization mechanism with
different numerical features of a piece.
He describes the KKPY test using a nonzero
Hochschild degree-two part and a Hodge
subspace of dimension at most two: a
projective surface contributing that part
would require at least three Hodge classes.
\item He mentions Gu\'er\'e's use of the
rank of the operator and Fay's finer
quadratic-form invariants. He does not
state the precise latter result, so no
precise theorem is inferred from the
speech-recognition transcript.
\end{itemize}

\medskip\noindent
\textbf{Bibliographical pointers.}
The preceding results are attributed to
the talks and their displayed references.
The following source identities and the
indicated result numbers were checked:
\begin{itemize}
\item T. Hinault, T. Y. Yu, C. Zhang,
S. Zhang,
\emph{Decomposition and framing of
$F$-bundles and applications to quantum
cohomology}, Proceedings of the London
Mathematical Society 133 (2026), e70209.
\url{https://doi.org/10.1112/plms.70209}.
Theorems 1.1 and 1.2 in arXiv version 2
are the formal and non-Archimedean
spectral decomposition statements:
\url{https://arxiv.org/abs/2411.02266v2}.
\item H. Iritani,
\emph{Quantum cohomology of blowups},
\url{https://arxiv.org/abs/2307.13555}.
\item H. Iritani,
\emph{Notes on the decomposition theorem
for blowups},
\url{https://arxiv.org/abs/2604.10028}.
\item V. Benedetti, A. Fay, J. Gu\'er\'e,
L. Manivel, N. Perrin,
\emph{An atomic criterion for irrationality
without quantum computations},
\url{https://arxiv.org/abs/2607.26718}.
Version 1: Proposition 2.2 and Corollary 2.4
(surface bounds), Proposition 3.1
(Hodge generality), Theorem 4.1
(the criterion), and Proposition 5.1
(the three families of examples).
\item V. Benedetti, L. Manivel, N. Perrin,
\emph{Quantum cohomology and irrationality
of Gushel-Mukai fourfolds},
\url{https://arxiv.org/abs/2603.17487}.
\end{itemize}
\end{document}
